%% confusable_sdp.tex
%%
%% A polynomial-optimisation approach to finding and certifying the
%% absence of confusable multisets for dotting encoders.
%%
%% Companion document to algorithm_description.tex.
%% Written with Claude's assistance, June 2026.

\documentclass[12pt,a4paper]{article}

\usepackage{amsmath,amssymb,amsthm}
\usepackage{hyperref}
\usepackage{geometry}
\usepackage{listings}
\usepackage{xcolor}
\geometry{margin=2.5cm}

\newtheorem{definition}{Definition}
\newtheorem{proposition}{Proposition}
\newtheorem{remark}{Remark}
\newtheorem{observation}{Observation}

\lstset{
  language=Python,
  basicstyle=\small\ttfamily,
  keywordstyle=\bfseries,
  commentstyle=\color{gray},
  breaklines=true,
  frame=single,
  aboveskip=6pt,
  belowskip=6pt,
}

\title{Minimal Confusable Multisets:\\
A Polynomial-Optimisation Perspective}

\author{Lester (problem), Claude (document)}
\date{June 2026}

\begin{document}
\maketitle

\begin{abstract}
This document takes a step back from the even-odd hypercube construction
used in \texttt{MinimalConfusableSets} and asks: how might one find the
\emph{globally} minimal confusable multiset pair for a given dotting
encoder, without any structural assumptions?  We reformulate the problem
as a polynomial feasibility question, derive an exact algebraic criterion
for the smallest non-trivial case ($s = 2$), introduce the
sum-of-squares~(SOS) / semidefinite programming~(SDP) framework for
certifying the absence of small confusable pairs, and exhibit an explicit
family of size-3 confusable pairs that arises from a cyclic algebraic
construction entirely unrelated to the hypercube.  Working code is
provided in \texttt{confusable\_search.py} and \texttt{confusable\_sos.py}.
\end{abstract}

\tableofcontents
\newpage

%% -----------------------------------------------------------------------
\section{Problem statement}
\label{sec:problem}

Fix $m$ encoding directions $\vec{d}_1, \ldots, \vec{d}_m \in \mathbb{R}^k$
(the \emph{good bats}).  A dotting encoder maps an $s$-element multiset
$X = \{\vec{v}_1, \ldots, \vec{v}_s\} \subset \mathbb{R}^k$ to the
$ms$-tuple consisting, for each $\vec{d}_j$, of the $n$ dot products
$\vec{v}_i \cdot \vec{d}_j$ sorted in non-decreasing order.

\begin{definition}
  Two distinct multisets $X, Y \subset \mathbb{R}^k$ of size $s$ are
  \emph{confusable} if, for every encoding direction $\vec{d}_j$, the
  sorted tuples of dot products with $\vec{d}_j$ are identical:
  \[
    \text{sort}\!\left(\{\,\vec{v} \cdot \vec{d}_j : \vec{v} \in X\,\}\right)
    =
    \text{sort}\!\left(\{\,\vec{v} \cdot \vec{d}_j : \vec{v} \in Y\,\}\right).
  \]
\end{definition}

The central question of this document is: for given $m$, $k$, and a
specific choice of encoding directions, what is the \emph{smallest} $s$
for which a confusable pair of size $s$ exists?  We make no assumption
about the form that confusable pairs take---in particular, we do not
assume they arise from the even-odd hypercube construction.

\subsection*{Power-sum reformulation}

Two multisets of real numbers are equal (as multisets) if and only if
all their power sums agree.  Applied to the scalar projections
$\{\vec{v} \cdot \vec{d}_j : \vec{v} \in X\}$, two size-$s$ multisets
are confusable if and only if

\begin{equation}
  \label{eq:powersums}
  \sum_{\vec{v} \in X} (\vec{v} \cdot \vec{d}_j)^r
  =
  \sum_{\vec{v} \in Y} (\vec{v} \cdot \vec{d}_j)^r
  \qquad \text{for all } j = 1, \ldots, m \text{ and } r = 1, \ldots, s.
\end{equation}

This is a system of $ms$ polynomial equations in the $2sk$ coordinates
of the vectors in $X$ and $Y$.  A confusable pair of size $s$ exists if
and only if this system has a solution with $X \neq Y$ as multisets.

\begin{remark}
  The minimum confusable-pair size is the smallest $s$ for which this
  polynomial system has a non-trivial solution.  It is a property of the
  specific set of directions $\{\vec{d}_j\}$, not of any particular
  construction method.
\end{remark}

%% -----------------------------------------------------------------------
\section{The case \texorpdfstring{$s = 2$}{s = 2}: exact characterisation}
\label{sec:s2}

For $s = 2$ the polynomial system~\eqref{eq:powersums} simplifies
substantially.  Write $X = \{\vec{x}_1, \vec{x}_2\}$ and
$Y = \{\vec{y}_1, \vec{y}_2\}$.  The $r = 1$ conditions force the
centroids to agree: $\vec{x}_1 + \vec{x}_2 = \vec{y}_1 + \vec{y}_2$
(when the $\vec{d}_j$ span $\mathbb{R}^k$).  Setting the common centroid
to zero (WLOG), write $\vec{x}_i = \pm\vec{p}/2$ and
$\vec{y}_i = \pm\vec{q}/2$.  The $r = 2$ conditions then reduce to

\[
  |\vec{p} \cdot \vec{d}_j| = |\vec{q} \cdot \vec{d}_j|
  \quad \forall j,
\]

i.e., for each direction $j$: either $(\vec{p} + \vec{q}) \cdot \vec{d}_j
= 0$ or $(\vec{p} - \vec{q}) \cdot \vec{d}_j = 0$ (or both).  Setting
$\vec{A} = \vec{p} + \vec{q}$ and $\vec{B} = \vec{p} - \vec{q}$:

\begin{proposition}
  \label{prop:s2}
  A confusable pair of size~$2$ exists if and only if there exist
  non-zero $\vec{A}, \vec{B} \in \mathbb{R}^k$ with $\vec{A} \neq \pm
  \vec{B}$ such that every encoding direction satisfies
  $\vec{d}_j \in \vec{A}^\perp \cup \vec{B}^\perp$.
\end{proposition}

In $\mathbb{R}^2$ this says: the $m$ encoding directions can be
partitioned into two groups, one whose members are all parallel to a
fixed direction $\vec{A}^\perp$, and one whose members are all parallel
to $\vec{B}^\perp$.  For generic directions (no two parallel), this
requires $m \leq 2$.

\subsection*{Worked example}

Consider $k = 2$ with two directions $\vec{d}_1 = (1, 0)$ (at $0°$) and
$\vec{d}_2 = (\tfrac{1}{\sqrt{2}}, \tfrac{1}{\sqrt{2}})$ (at $45°$).
Take $\vec{A} \propto (0, 1)$ (so $\vec{A}^\perp$ is the horizontal
axis, containing $\vec{d}_1$) and $\vec{B} \propto (1, -1)$ (so
$\vec{B}^\perp$ is the $45°$ axis, containing $\vec{d}_2$).  These are
non-zero and non-parallel.  Then

\[
  X = \left\{\tfrac{1}{4}(1, 1),\; -\tfrac{1}{4}(1, 1)\right\},
  \qquad
  Y = \left\{\tfrac{1}{4}(1, -1),\; -\tfrac{1}{4}(1, -1)\right\}
\]

is a confusable pair of size~2.  One may verify directly:
both $X$ and $Y$ project to the sorted tuple $(-\tfrac{1}{4},
\tfrac{1}{4})$ onto $\vec{d}_1$, and to $(-\tfrac{1}{2\sqrt{2}},
\tfrac{1}{2\sqrt{2}})$ onto $\vec{d}_2$.

Adding any third generic direction immediately destroys this: a
direction at $\theta \neq 0°, 45°$ does not lie in
$\vec{A}^\perp \cup \vec{B}^\perp$, so Proposition~\ref{prop:s2} fails,
and the SOS certificate in Section~\ref{sec:sos} confirms injectivity.

%% -----------------------------------------------------------------------
\section{An explicit \texorpdfstring{$s = 3$}{s = 3} confusable family}
\label{sec:cyclic}

A confusable pair of size~3 for the three directions
$\vec{d}_1 = (1, 0)$, $\vec{d}_2 = (\tfrac{1}{\sqrt{2}},
\tfrac{1}{\sqrt{2}})$, $\vec{d}_3 = (0, 1)$ can be constructed
\emph{without} any reference to the even-odd hypercube.  For any triple
of distinct real numbers $(a, b, c)$, define

\[
  X = \{(a, b),\; (c, a),\; (b, c)\},
  \qquad
  Y = \{(a, c),\; (b, a),\; (c, b)\}.
\]

\begin{proposition}
  \label{prop:cyclic}
  For any $(a, b, c)$ with $\{a, b, c\}$ not all equal, the pair
  $(X, Y)$ defined above is confusable with respect to all three
  directions $\vec{d}_1$, $\vec{d}_2$, $\vec{d}_3$.
\end{proposition}

\begin{proof}
  For $\vec{d}_1 = (1, 0)$: both $X$ and $Y$ have first-coordinates
  $\{a, c, b\}$ as a multiset.

  For $\vec{d}_3 = (0, 1)$: both have second-coordinates $\{b, a, c\}$.

  For $\vec{d}_2 = (\tfrac{1}{\sqrt{2}}, \tfrac{1}{\sqrt{2}})$: the
  diagonal projections of $X$ are proportional to $\{a+b,\, c+a,\,
  b+c\}$ and those of $Y$ to $\{a+c,\, b+a,\, c+b\}$.  Since addition
  is commutative, both sets are $\{a+b,\, a+c,\, b+c\}$.

  In each case the sorted projections are identical.  Since $X$ and $Y$
  differ as multisets when $a, b, c$ are not all equal (the first
  coordinates are the same set but paired with different second
  coordinates), $(X, Y)$ is a genuine confusable pair.
\end{proof}

\begin{remark}
  This construction is entirely unrelated to the even-odd hypercube: the
  elements of $X$ and $Y$ are not subset sums of any collection of bad
  bats.  It illustrates that the true minimum confusable-pair size may be
  achieved by structures outside the hypercube framework.
\end{remark}

For a concrete numerical instance: $(a, b, c) = (5.5, -5.7, 0.2)$ gives
\[
  X = \{(5.5, -5.7),\; (0.2, 5.5),\; (-5.7, 0.2)\},
  \quad
  Y = \{(5.5, 0.2),\; (-5.7, 5.5),\; (0.2, -5.7)\},
\]
a confusable pair of size~3, confirmed by the numerical code.

\subsection*{When does a fourth direction destroy confusability?}

A direction $\vec{d} = (\cos\theta, \sin\theta)$ with $\theta \neq 0°,
45°, 90°$ will generically break the size-3 confusable pair: the cubic
power sum (the $r = 3$ condition) for direction $\vec{d}$ is

\[
  p_3^X(\theta) - p_3^Y(\theta)
  = 3\cos^2\!\theta\sin\theta\,(a^2b + c^2a + b^2c - a^2c - b^2a - c^2b)
  + 3\cos\theta\sin^2\!\theta\,(ab^2 + ca^2 + bc^2 - ac^2 - ba^2 - cb^2).
\]

This vanishes for all $(a, b, c)$ only when $\cos^2\theta\sin\theta =
\cos\theta\sin^2\theta$, i.e.\ $\theta \in \{0°, 45°, 90°\}$.  Adding any
other direction $\vec{d}_4$ makes the $r = 3$ condition non-trivially
constraining, and numerical search (Section~\ref{sec:numerical}) finds no
confusable pairs of size up to~5 for four equally-spaced directions.

%% -----------------------------------------------------------------------
\section{Sum-of-squares certificates for size-2 injectivity}
\label{sec:sos}

\subsection*{The certificate polynomial}

For the $s = 2$ question, define
\[
  f(\vec{A}, \vec{B})
  = \sum_{j=1}^m (\vec{A} \cdot \vec{d}_j)^2\,(\vec{B} \cdot \vec{d}_j)^2,
  \qquad
  g(\vec{A}, \vec{B}) = \|\vec{A}\|^2\,\|\vec{B}\|^2.
\]

From Proposition~\ref{prop:s2}, a confusable pair of size~2 exists if
and only if $f$ has a zero for some non-zero $(\vec{A}, \vec{B})$.  This
is equivalent to $f$ failing to satisfy $f \geq \lambda g$ for any
$\lambda > 0$.  Conversely:

\begin{observation}
  If there exists $\lambda > 0$ such that $f - \lambda g$ is a
  sum-of-squares (SOS) polynomial, then $f(\vec{A}, \vec{B}) \geq
  \lambda \|\vec{A}\|^2 \|\vec{B}\|^2 > 0$ for all $(\vec{A}, \vec{B})
  \neq \vec{0}$, and hence no confusable pair of size~2 exists.
\end{observation}

Note that $f$ itself is already SOS ($f = \sum_j [(\vec{A} \cdot
\vec{d}_j)(\vec{B} \cdot \vec{d}_j)]^2$), so $\lambda = 0$ is always
feasible.  The question is whether the optimal $\lambda^* > 0$.

\subsection*{SDP formulation}

An SOS polynomial of degree~4 in $2k$ variables can be written as
$v^T Q v$ where $v$ is the vector of all monomials of degree $\leq 2$ in
the $2k$ variables $(\vec{A}, \vec{B})$ and $Q$ is a PSD (positive
semidefinite) matrix.  The number of such monomials is $\binom{2k +
2}{2}$, giving a $\binom{2k+2}{2} \times \binom{2k+2}{2}$ matrix.  For
$k = 2$: a $15 \times 15$ matrix.

The SDP is:
\begin{align*}
  \text{maximise} \quad & \lambda \\
  \text{subject to} \quad
    & Q \succeq 0, \\
    & v^T Q v = f - \lambda g \quad\text{(as polynomial identity).}
\end{align*}

The polynomial identity is enforced by matching the coefficient of each
degree-4 monomial $z^\gamma$ on both sides:
\[
  \sum_{\substack{\alpha + \beta = \gamma \\ |\alpha|, |\beta| \leq 2}}
  Q_{\alpha\beta}
  \;=\;
  f_\gamma - \lambda\, g_\gamma,
\]
where $f_\gamma$ and $g_\gamma$ are the known coefficients of $f$ and
$g$ at monomial $\gamma$.  These are linear equalities in the decision
variables $(Q, \lambda)$, so the overall problem is a standard SDP.

\subsection*{Results for $k = 2$}

Running the SDP (implemented in \texttt{confusable\_sos.py}) for three
choices of encoding directions in $\mathbb{R}^2$:

\medskip
\begin{center}
\begin{tabular}{lll}
  \hline
  Directions & $\lambda^*$ & Conclusion \\
  \hline
  $\vec{d}_1 = (1,0)$,\; $\vec{d}_2 = (\frac{1}{\sqrt2}, \frac{1}{\sqrt2})$ (2 dirs)
    & $0.000$ & Confusable size-2 pair may exist \\
  Add $\vec{d}_3 = (0,1)$ (3 dirs)
    & $0.167$ & \textbf{Certified}: no confusable size-2 pair \\
  Add $\vec{d}_4 = (-\frac{1}{\sqrt2}, \frac{1}{\sqrt2})$ (4 dirs)
    & $0.500$ & \textbf{Certified}: no confusable size-2 pair \\
  \hline
\end{tabular}
\end{center}
\medskip

For 2 directions, $\lambda^* = 0$ confirms that a confusable pair
exists (as constructed in Section~\ref{sec:s2}).  For 3 and 4 directions,
$\lambda^* > 0$ provides a rigorous certificate of injectivity for
size-2 multisets.

%% -----------------------------------------------------------------------
\section{Numerical search for larger \texorpdfstring{$s$}{s}}
\label{sec:numerical}

For $s \geq 3$, the polynomial system~\eqref{eq:powersums} is harder to
analyse algebraically, and the SDP certificate needs to be extended to
higher-dimensional moment matrices (see Section~\ref{sec:lasserre}).
A practical complement is a gradient-based numerical search.

\subsection*{Search strategy}

For fixed $s$, $k$, and directions $D$, the code in
\texttt{confusable\_search.py} proceeds as follows:

\begin{enumerate}
  \item Sample $X \in \mathbb{R}^{s \times k}$ at random with zero
    centroid.
  \item Compute the target moments $t = (p_{r,j}(X))_{r,j}$ via
    equation~\eqref{eq:powersums}.
  \item From a \emph{random} starting point $Y_0$ (not near $X$),
    minimise the confusability loss $\|p(Y) - t\|^2$ over $Y$ using
    L-BFGS-B.
  \item If the minimised loss falls below $10^{-10}$ and the Hungarian
    matching distance between $X$ and $Y$ exceeds $10^{-2}$, report a
    confusable pair.
\end{enumerate}

The random initialisation of $Y_0$ is critical: starting near $X$
causes the optimiser to converge trivially to $Y = X$.

\subsection*{Results for $k = 2$}

Running 500 random trials for each $(s, m)$ pair:

\medskip
\begin{center}
\begin{tabular}{cclll}
  \hline
  $s$ & $m = 2$ dirs & $m = 3$ dirs & $m = 4$ dirs \\
  \hline
  2 & \textbf{Found} & not found & not found \\
  3 & \textbf{Found} & \textbf{Found} & not found \\
  4 & \textbf{Found} & \textbf{Found} & not found \\
  5 & \textbf{Found} & \textbf{Found} & not found \\
  \hline
\end{tabular}
\end{center}
\medskip

With only 2 directions confusable pairs are found for every $s \geq 2$,
as expected (two directions generically leave a 1-dimensional family of
confusable pairs at each size).  With 3 directions confusable pairs exist
for $s \geq 3$ (including the cyclic family of
Section~\ref{sec:cyclic}).  With 4 equally-spaced directions, no
confusable pair was found for any $s \leq 5$ in 500 trials, which gives
numerical evidence in support of the conjecture that 4 directions suffice
to inject all multisets of size $\leq 5$ of 2-vectors.

\begin{remark}
  Failure to find a confusable pair by this search is \emph{not} a proof
  of injectivity: it may simply mean that confusable pairs exist but are
  rare or have small basins of attraction.  A rigorous lower bound
  requires the Lasserre hierarchy (Section~\ref{sec:lasserre}).
\end{remark}

%% -----------------------------------------------------------------------
\section{The Lasserre hierarchy for general \texorpdfstring{$s$}{s}}
\label{sec:lasserre}

\subsection*{Moment relaxation}

The existence of a confusable pair of size~$s$ is equivalent to the
feasibility of the polynomial system~\eqref{eq:powersums} with
$X \neq Y$.  The \emph{Lasserre moment hierarchy} provides a sequence of
SDPs, parametrised by a \emph{relaxation level} $d = 1, 2, \ldots$, that
gives increasingly tight approximations to this feasibility question.

At level $d$, the decision variables are the \emph{moment variables}
$y_\alpha$ for all multi-indices $\alpha$ with $|\alpha| \leq 2d$, where
$y_\alpha$ represents the expectation $\mathbb{E}[z^\alpha]$ of monomial
$z^\alpha$ under a hypothetical probability measure supported on the
feasible set.  The SDP constraints are:
\begin{itemize}
  \item \emph{Moment matrix}: $M_d(y) \succeq 0$ (a $\binom{n+d}{d}
    \times \binom{n+d}{d}$ PSD matrix, where $n = 2sk$ is the number of
    scalar decision variables in $X$ and $Y$).
  \item \emph{Constraint moments}: $\sum_\alpha c_\alpha y_\alpha = 0$
    for each equality constraint in \eqref{eq:powersums} (linear in $y$).
\end{itemize}
If this SDP is infeasible at level $d$, then no confusable pair of
size~$s$ exists.  If it is feasible, the solution can sometimes be
\emph{flat-extended} to recover an explicit confusable pair.

\subsection*{Complexity}

The moment matrix at level $d$ has $\binom{n+d}{d}$ rows and columns.
For $k = 2$, $s = 4$: $n = 16$, and at $d = 2$ the moment matrix is
$\binom{18}{2} = 153 \times 153$.  This is tractable with a modern SDP
solver (MOSEK, SCS, or Clarabel).  For $s = 5$: $n = 20$, moment matrix
$\binom{22}{2} = 231 \times 231$ at $d = 2$.

A practical implementation in Python would use \texttt{cvxpy} to
construct the moment matrix and add the confusability constraints as
linear equality constraints on the moments.  The
\texttt{confusable\_sos.py} file implements the $s = 2$ case directly;
extending it to general $s$ follows the same pattern with a larger
moment matrix.

%% -----------------------------------------------------------------------
\section{Code notes}
\label{sec:code}

Two companion Python files are provided, both requiring \texttt{numpy},
\texttt{scipy}, \texttt{sympy}, and \texttt{cvxpy} (\texttt{pip install
cvxpy sympy}).

\subsection*{\texttt{confusable\_search.py}}

Implements the gradient-based search described in
Section~\ref{sec:numerical}.  Main entry point:

\begin{lstlisting}
result = search_confusable_pair(s, k, directions,
                                n_restarts=500, res_tol=1e-10)
# Returns (X, Y, residual) if found, None otherwise.
\end{lstlisting}

The function accepts any $s$, $k$, and list of direction vectors.
Running the script directly exercises the cases in the table above.

\subsection*{\texttt{confusable\_sos.py}}

Implements:
\begin{itemize}
  \item \texttt{sos\_lower\_bound(directions)}: the SDP from
    Section~\ref{sec:sos}, maximising $\lambda$ s.t.\ $f - \lambda g$
    is SOS.  Returns $\lambda^*$.  If $\lambda^* > 0$, no confusable
    size-2 pair exists.
  \item \texttt{exact\_check\_s2(directions)}: an exhaustive $O(2^m)$
    combinatorial check of Proposition~\ref{prop:s2} for small $m$.
    Returns \texttt{(True, A, B, X, Y)} if a confusable pair is found.
\end{itemize}

\begin{lstlisting}
lam = sos_lower_bound(directions)  # SDP certificate
found, A, B, X, Y = exact_check_s2(directions)  # Combinatorial check
\end{lstlisting}

\end{document}
